% !TeX root = ../../main.tex
\subsection{Local transport and common dented regions}
\label{subsec:b:local-transport}

We next integrate the endpoint bounds along a characteristic. The
small part of this integral gains half a power of $S$. Its complement
is controlled by the order of vanishing at $y=0$, and then by the
regular-disk estimate. These three bounds are needed together; a
fixed-characteristic formal identity alone would not give convergence
on a common complex region.

\begin{lmm}[The local transport bound]\label{lmm:b:local-transport}
Let $\mathcal H$ have scale-$a$ endpoint bounds, $a\geq0$, and
let $\mathcal Q$ have the same bounds with zero square-root
coefficients. Suppose both vanish to an integer order
$N>a+1/2$ at $y=0$. Assume their bounds on a parameter interval
containing $[\theta,\theta+\sigma_0]$, where $\sigma_0>0$ is fixed
and sufficiently small. Then
\begin{equation}\label{eq:b:local-transport}
 \mathcal T(\theta;x,y)
 =\int_0^{\sigma_0}
       [(1+y\partial_y)\mathcal H+\mathcal Q]
       (\theta+\sigma;x+\sigma v_\theta(x),y)\,\dd\sigma
\end{equation}
has scale-$(a+1/2)$ endpoint bounds on every strictly smaller cone,
after reducing the small-$|S|$ range. Its coefficient of
$\sqrt{S\mp y}$ is
\begin{equation}\label{eq:b:transport-endpoint-coefficient}
 \frac{S}{x_{*,\theta}-v_\theta(x)}H_1^\pm(\theta,S).
\end{equation}
The value remainder has exponent $\nu$, and hence on smaller
domains obeys \eqref{eq:b:remainder-derivatives} with any fixed
derivative order, in particular orders one and two. Constants depend
only on the input bounds, $a,N,\nu$, the angular margins and the
fixed parameter intervals.
\end{lmm}
\begin{prf}
We give the positive-endpoint argument; the other sign is identical.
Write $d=x_{*,\theta}-v_\theta(x)$ locally, so
$|d|\asymp\sqrt{|S|}$ and $\operatorname{Re}d\geq c\sqrt{|S|}$.
For $\sigma$ at most a small multiple of $\sqrt{|S|}$,
Lemma~\ref{lmm:b:cone-transport} keeps the arguments in the larger
input cone and gives
\[
 |S_\sigma-y|\geq c\bigl(|S-y|+\sigma\sqrt{|S|}\bigr),
 \qquad S_\sigma-y=(S-y)+\sigma d+O(\sigma^2).
\]
The scale of the integrand is $|S|^a$, times an inverse square root
in the normalized distance, or a less singular remainder.
Its integral is bounded by $C|S|^{a+1/2}$.

\smallskip\noindent
\emph{The frozen singular term.}
The only nonintegrated square-root term differentiated by
$y\partial_y$ is $-yH_1^+(\theta,S)/(2\sqrt{S-y})$.
Freeze its numerator at $\sigma=0,y=S$ and its distance at
$(S-y)+\sigma d$. The model is
\begin{equation}\label{eq:b:frozen-model}
 -\frac{S H_1^+(\theta,S)}2
       \int_0^{c\sqrt{|S|}}
                  ((S-y)+\sigma d)^{-1/2}\,\dd\sigma.
\end{equation}
The upper-end value of its elementary primitive is analytic in $y$
near $S$. Its lower-end value is exactly
\eqref{eq:b:transport-endpoint-coefficient} times $\sqrt{S-y}$.
No coefficient is divided out in this calculation, so it also applies
when $H_1^+$ vanishes.

Here are uniform error bounds for these replacements. Put
$u=\sigma/\sqrt{|S|}$ only in this estimate. Cauchy's bounds for
the numerator give the additive error
\[
 |yH_1^+(\theta+\sigma,S_\sigma)-S H_1^+(\theta,S)|
       \leq C|S|^{a+1/2}(u+|1-y/S|).
\]
The distance error is $O(|S|u^2)$. After one normalized
$y$ derivative, each resulting error density is bounded, apart
from its factor $C|S|^a$, by
$(|1-y/S|+u)^{-1/2}$.
For the distance error this follows from
$u^2(|1-y/S|+u)^{-5/2}\leq(|1-y/S|+u)^{-1/2}$;
the numerator error is bounded in the same way.
These bounds are integrable in $u$. After subtracting their endpoint
values the errors are $O(|S|^{a+1/2}|1-y/S|)$.

For the input remainder, the normalized derivative of its Euler term
has majorant
$C|S|^a(|1-y/S|+u)^{\nu-2}$.
Its integral is $O(|S|^{a+1/2}|1-y/S|^{\nu-1})$.
Integration in the endpoint variable gives a value remainder
$O(|S|^{a+1/2}|1-y/S|^\nu)$.
The undifferentiated square-root term and $\mathcal Q$ have bounded
first normalized derivative after the $u$ integration, and therefore
give only a constant and a remainder of this size.

\smallskip\noindent
\emph{The rest of the local characteristic.}
Between $c\sqrt{|S|}$ and $C\sqrt{|S|}$, with fixed positive
$c,C$, the endpoint distances are bounded below by a multiple of
$|S|$. This part is analytic near $y=S$ and has norm
$O(|S|^{a+1/2})$. Choose $C$ large. Beyond $C\sqrt{|S|}$ the
ratio $|y/S_\sigma|$ is uniformly less than one, by
\eqref{eq:b:transport-distances}. The vanishing order $N$ and
Cauchy's estimate on a larger inner disk give the bound
\begin{equation}\label{eq:b:intermediate-transport}
 C|S|^N\int_{C\sqrt{|S|}}^{\sigma_0}
                    (\sqrt{|S|}+\sigma)^{2a-2N}\,\dd\sigma
                 \leq C|S|^{a+1/2}.
\end{equation}
The integral converges uniformly at its scaled upper end because
$N>a+1/2$. The same estimate on slightly larger endpoint
neighborhoods shows that this part contributes an analytic constant
and an $O(|S|^{a+1/2}|1-y/S|)$ variation.
This argument does not assume that $|S_\sigma|$ increases along a
complex characteristic; only the two-distance bounds are used.

Together these estimates prove the value expansion and coefficient
claimed. Away from the endpoints the normalized arguments range over
a compact set; the same three characteristic estimates apply with
their endpoint factors omitted. Normal convergence makes its endpoint
coefficients analytic
in the remaining parameters. On smaller angular domains Cauchy's
formula regenerates every fixed derivative bound from the value
remainder, including the third order needed in a subsequent layer.
This completes the proof.
\end{prf}

The next proposition records analytic existence and uniform bounds
for the actual fields, without identifying their leading coefficients.
That identification and its iterated integral are used in Section 8.
For real parameters $r>0$, $0<\gamma<\pi/4$ and $S>0$,
a doubly dented disk below means the open set
\begin{equation}\label{eq:b:dented-disk}
 |y|<(1+r)S,\qquad
 |\arg(S-y)|<\pi-\gamma,\qquad
 |\arg(S+y)|<\pi-\gamma.
\end{equation}
The two endpoints themselves are omitted. For a sufficiently small
$r>0$, this set lies in a smaller two-distance cone than
the one used for its estimates. To see this, use compactness on
$|y|\leq S$ away from the endpoints; near an endpoint one distance
has argument close to zero, and the other has the displayed fixed
margin from the negative axis. Both observations persist after a
small enlargement of the normalized radius.

\begin{prp}[Common Borel domains]\label{prp:b:common-domains}
For every fixed integer $H\geq1$ and $1/2<\nu<1$ there are
$\theta_H,r_H>0$ and $0<\gamma_H<\pi/4$ such that, for $1\leq h\leq H$,
$0\leq\theta\leq\theta_H$ and real
$-r_H\leq x<R_\theta$, the actual Borel series
$\mathcal L_h(\theta;x,y)$ converges for $|y|<S=R_\theta-x$.
It extends holomorphically to \eqref{eq:b:dented-disk} with
$r=r_H,\gamma=\gamma_H$. These extensions are joint
in $x,y$ on complex neighborhoods of the indicated real centers.
Near the critical center they include the smaller two-distance
cones above, with all constants uniform. One has
\begin{equation}\label{eq:b:common-bounds}
 \begin{aligned}
 |\mathcal L_h|&\leq C_H|S|^{(h-1)/2},\\
 |\partial_x^a\partial_y^k\mathcal L_h|
 &\leq C_{H,a,k}|S|^{(h-1)/2-a-k}
   \left(1+|1-y/S|^{1/2-a-k}
            +|1+y/S|^{1/2-a-k}\right)
                  \quad(a+k\geq1).
 \end{aligned}
\end{equation}
Here $a,k$ are fixed nonnegative integers and derivative estimates
use smaller angular domains. The finite-order remainder class
\eqref{eq:b:endpoint-class} holds with scale $(h-1)/2$ and exponent
$\nu$. All kernel sums and characteristic integrals converge
normally on compact subsets of these regions.
\end{prp}
\begin{prf}
For $h=1$, insert the local square-root inverse
\eqref{eq:b:morse-inverse} into the exact log quotient. Its leading
local expression, up to bounded analytic terms, is
\[
 \log\frac{2\sqrt S}{\sqrt{S-y}+\sqrt S}.
\]
The odd difference removes every term independent of $y$.
It is uniformly bounded on the two-distance cones, with a
square-root endpoint expansion and scale zero. The $m\geq2$
terms are analytic at the two endpoints and are normally summable:
their odd differences have the additional $m^{-1}$ factor.
The same argument at a real center away from the critical point
uses an analytic nonzero difference quotient and the local
square-root expansion of its endpoint inverse. Its constants are
uniform on compact center intervals. This proves the initial
endpoint class as well as its local holomorphic extension.

Induct simultaneously on normal convergence and this endpoint class.
If the lower-layer scales are $(i-1)/2$,
Lemmas~\ref{lmm:b:linear-endpoint} and
\ref{lmm:b:quadratic-endpoint} give scale $(h-2)/2$ for the two
sources. Their exact order at $y=0$ is at least $N=3h-2$, so
$N>(h-2)/2+1/2$. Lemma~\ref{lmm:b:local-transport} applies to
the initial characteristic segment. Its complement is the normally
convergent remote integral of Corollary~\ref{crl:b:remote}; near
the critical center it is $O(|y|^N)$ and has no endpoint singularity.
We obtain the scale $(h-1)/2$ and the claimed remainder class.
The resulting functions agree with the formal coefficients on a
small inverse disk, by normal coefficient extraction, and hence
with the actual series already constructed in
Proposition~\ref{prp:b:tube}.

The choices can be made in a finite order: use successively smaller
cones, $S$ neighborhoods and $\theta$ intervals for increasing $h$.
At a new layer choose $\sigma_0$ so small that the lower-layer
parameter interval contains $[\theta,\theta+\sigma_0]$ and that
the moving critical points stay in their Morse neighborhoods.
The remote estimate needs no critical point beyond this interval.
Choose the new $\theta$ interval inside half the preceding one.
Since $H$ is fixed, all resulting constants are positive.
Value remainders on the larger cone give every fixed derivative
order on a smaller one, so this construction does not consume a
finite stock of derivative estimates.

\smallskip\noindent
\emph{Patching the common complex region.}
For real centers bounded away from the critical point, the inverse
radius and the noncritical-circle continuation are provided by
Propositions~\ref{prp:b:inverse-budget} and \ref{prp:b:tube}.
Their compact normalized arcs have common neighborhoods.
Near $y=S$, the principal shifted center has
$S'=\tau S+(1-\tau)(S-y)$ and its new endpoint distance is
$S-y$. If $\tau$ is at most a fixed multiple of $|1-y/S|$,
both belong to the small critical cone already constructed.
Otherwise the center is a small complex perturbation of the real
segment, and its normalized endpoint distance is smaller by a
factor comparable to $\tau$. This is precisely the split used in
Lemma~\ref{lmm:b:linear-endpoint}; it supplies the joint endpoint
neighborhoods with uniform bounds. The negative branch is identical.
All other shifted arguments have a strict budget except at lower-layer
endpoints, and the quadratic endpoint corners are covered by
Lemma~\ref{lmm:b:quadratic-endpoint}. A finite cover joins these
endpoint patches to the noncritical arcs. Small negative real centers
are allowed by the negative extension of the inverse-radius budget.
Finally take the minimum of these finitely many normalized radius
and angle margins and the small-critical cone margins.
This gives \eqref{eq:b:dented-disk} uniformly, including centers
tending to the critical point.

The same constructions are holomorphic in $\theta$ near its compact
real interval. In the critical patches this follows directly from
the analytic Morse coordinate: $x_{*,\theta}'$ remains close to a
positive real number, so the two-distance estimates retain their
angular margins for a small complex parameter perturbation. The
remaining compact inverse patches and the regular remote bounds
persist under that perturbation as well. Thus Cauchy's formula
justifies the parameter and coefficient derivatives used in the
induction. Formula~\eqref{eq:b:input-derivatives} gives
\eqref{eq:b:common-bounds}. This completes the analytic construction
without assuming convergence or continuation of a higher Borel layer.
\end{prf}
